%! TEX root = main.tex

\subsection{Roadmap: from intuition to proofs}

You already know how to compute many limits and integrals.  This course is
about stating precisely what those computations mean, and proving that the
rules you use are valid.  We will proceed in this order:
\begin{enumerate}
  \item limits and their $\varepsilon$--$\delta$ definition;
  \item continuity and its consequences (especially the intermediate value
    and extreme value theorems);
  \item differentiability, the mean value theorem, and rigorous local
    approximation;
  \item sequences and series, including convergence tests and uniform ideas
    where they matter;
  \item the Riemann integral and the theorem connecting integration with
    antiderivatives.
\end{enumerate}
The dependency graph is real: the completeness of $\R$ drives the main global
theorems about continuous functions, while sequences give a second, often more
usable, language for limits.

\subsection{Limits: the local contract}

The informal statement ``$f(x)$ gets close to $L$ when $x$ gets close to
$a$'' hides a contract between two players.  An adversary first chooses any
desired output accuracy $\varepsilon>0$.  You must then give an input radius
$\delta>0$ which guarantees that accuracy.  Crucially, your $\delta$ may
depend on $\varepsilon$, but not on the particular $x$ later chosen inside
that radius.

\begin{definition}[Limit at a point]
Let $D\subseteq\R$, let $f:D\to\R$, and suppose that $a$ is an
\emph{accumulation point} of $D$: every interval around $a$ contains a point
of $D$ distinct from $a$.  We write
\[
  \lim_{x\to a}f(x)=L
\]
if, for every $\varepsilon>0$, there is a $\delta>0$ such that for every
$x\in D$,
\[
  0<\abs{x-a}<\delta \quad\Longrightarrow\quad \abs{f(x)-L}<\varepsilon.
\]
\end{definition}

The exclusion $x\ne a$ is not decoration: a limit describes values nearby,
not necessarily at, $a$.  Thus changing $f(a)$ never changes a limit at $a$.
The accumulation-point condition prevents a vacuous definition---without
nearby domain points, any $L$ would satisfy the implication.

\begin{example}
Prove $\lim_{x\to 3}(2x-1)=5$.  Let $\varepsilon>0$.  Choose
$\delta=\varepsilon/2$.  If $0<\abs{x-3}<\delta$, then
\[
  \abs{(2x-1)-5}=2\abs{x-3}<2\delta=\varepsilon.
\]
Since this works for every $\varepsilon>0$, the limit is $5$.
\end{example}

This proof is the ideal simple shape: begin with an arbitrary $\varepsilon$,
announce a positive $\delta$, and turn the assumed bound on $\abs{x-a}$ into
the required bound.  Working backwards before writing the proof often reveals
the right choice of $\delta$.

\begin{example}
Prove $\lim_{x\to 2}x^2=4$.  We have
\[
  \abs{x^2-4}=\abs{x-2}\abs{x+2}.
\]
The second factor is troublesome: it is not controlled directly by the output
tolerance.  First impose the auxiliary condition $\abs{x-2}<1$.  Then
$1<x<3$, so $\abs{x+2}<5$.  Given $\varepsilon>0$, set
\[
  \delta=\min\left\{1,\frac{\varepsilon}{5}\right\}.
\]
For $0<\abs{x-2}<\delta$, the auxiliary bound gives
\[
  \abs{x^2-4}<5\abs{x-2}<5\delta\leq\varepsilon.
\]
\end{example}

The $\min$ is the standard way to satisfy multiple local requirements at
once.  It is analogous to establishing an invariant (here, $\abs{x+2}<5$)
before using it.  Unlike a program invariant, however, the claim quantifies
over a continuum of possible inputs, so the proof must establish the bound
symbolically.

\begin{theorem}[Uniqueness of limits]
If $\lim_{x\to a}f(x)=L$ and $\lim_{x\to a}f(x)=M$, then $L=M$.
\end{theorem}

\begin{proof}
Suppose instead that $L\ne M$, and put
$\varepsilon=\abs{L-M}/3>0$.  Obtain positive radii $\delta_L,\delta_M$
from the two limit statements for this same $\varepsilon$.  Since $a$ is an
accumulation point, choose $x\in D$ with
$0<\abs{x-a}<\min\{\delta_L,\delta_M\}$.  The triangle inequality gives
\[
  \abs{L-M}\leq \abs{L-f(x)}+\abs{f(x)-M}
  <2\varepsilon=\frac{2}{3}\abs{L-M},
\]
which is impossible.  Therefore $L=M$.
\end{proof}

\begin{remark}[Common proof failures]
\begin{itemize}
  \item Choosing $\delta$ after seeing $x$ reverses the quantifiers and is
  not a limit proof.
  \item Saying ``take $x$ sufficiently close'' is the conclusion to prove;
  name an explicit $\delta(\varepsilon)$.
  \item A bound such as $\abs{x+2}<5$ needs a stated reason, typically an
  auxiliary restriction like $\delta\leq1$.
  \item To disprove a proposed limit, negate the definition correctly:
  find some $\varepsilon_0>0$ such that every $\delta>0$ admits an $x$ with
  $0<\abs{x-a}<\delta$ and $\abs{f(x)-L}\geq\varepsilon_0$.
\end{itemize}
\end{remark}

